\begin{proof}

\begin{enumerate}[label={$\langle1\rangle$\arabic*.},leftmargin=*]
\item Let $\rho_N:T\to N$ be the minimal normalization resolution,
let $\nu:N\to S$ be the finite normalization, and put
$\rho=\nu\rho_N$. The reduced-horizontal case is already
proved by P4R.9. We therefore use the actual double-horizontal
pole divisor $D$ of P4R.6. Its support is a smooth section $H$
with rational trees attached. The pushout and curve-group
calculation of P4R.7--P4R.9 show that the reduced image
conductor $C$ is connected, that its point-identification
graph has no cycles, and that it has exactly one normalized
component of positive genus, of genus one. Denote the
corresponding reduced component by $C_E$ and its normalization
by $j:E\to C_E\hookrightarrow Y$. In particular every point
of $C_E$ is unibranch. Every other normalized component of
$C$ is rational.

\item A rational component of $D$ cannot map nonconstantly to $E$.
Thus $H$ maps nonconstantly to $E$. Write $g=g(H)\ge1$.
At a general point of $C_E$, every normalization branch has
conductor exponent two, because the coefficient of $H$ in
$D$ is two. The intrinsic conductor exponent of a singular
plane branch is at least two, and each additional branch
increases it by its positive intersection multiplicity.
The possibilities are therefore a single ordinary cusp,
two smooth branches with intersection multiplicity two, or
three smooth branches with pairwise intersection multiplicity
one. These are $A_2,A_3,D_4$, respectively, and $H\to E$
has degree $1,2,3$ in the three cases.

\item We first construct a map which detects global functions on
the entire conductor, including its rational components.
Since $H$ is not contracted, P4.33 proves that $N$ has
rational singularities. Put $L_N=\nu^*\omega_S$.
Finite duality identifies its conductor ideal with
\[
 I_N=\omega_N\otimes L_N^{-1}.
\]
The canonical trace is an isomorphism
$\rho_{N*}\omega_T=\omega_N$ for rational analytic
singularities. We use Greb, Remark~2.1, p.~415
\cite{Greb2011}. The hypotheses hold here: $N$ is normal
and Cohen--Macaulay, $\rho_N$ is a proper resolution,
P4.33 gives $R^1\rho_{N*}\mathcal O_T=0$, and its
one-dimensional fibres give vanishing of the higher
structure-sheaf direct images. The integral pole identity
of P4R.6 and projection formula consequently give
\[
 \rho_{N*}\mathcal O_T(-D)
 =\rho_{N*}(\omega_T\otimes\rho_N^*L_N^{-1})
 =I_N.
\]
These are equalities of ideal subsheaves of $\mathcal O_N$:
the pole identity uses the pullback of the same residue
section which defines finite duality, and the trace agrees
with it on the birational isomorphism locus. Equality there
determines the maps between these torsionfree sheaves.
The divisor ideal sequence therefore gives an injection
\[
 \mathcal O_N/I_N\longrightarrow\rho_{N*}\mathcal O_D.
\]
The finite conductor square gives an injection from the
image conductor scheme. Composing them yields the actual
morphism of coherent $\mathcal O_S$-algebras
\[
 \mathcal O_{C_{\rm sch}}\longrightarrow
 \nu_*(\mathcal O_N/I_N)\longrightarrow\rho_*\mathcal O_D,
 \qquad\hbox{both arrows injective}.                         \tag{P4R.13.1}
\]

\item We show that a global function on $D$ is determined by
its value at the generic doubled horizontal section.
At every point blowup $\beta:T'\to T''$ in the actual
map to the minimal ruled surface, the pole divisors obey
$D'=\beta^*D''-E_\beta$. The elementary direct images are
\[
 \beta_*\mathcal O(E_\beta)=\mathcal O,
 \quad R^1\beta_*\mathcal O(E_\beta)=0,
 \quad\beta_*\mathcal O=\mathcal O,
 \quad R^1\beta_*\mathcal O=0.
\]
For the first equality, a meromorphic function allowed
a pole only on $E_\beta$ descends off the centre and
extends across that point by Hartogs. For the vanishing,
the two blowup charts, or their exceptional-curve
filtration, have successive terms
$\mathcal O_{\mathbb P^1}(j)$ with $j\ge-1$;
their first cohomology vanishes. Equivalently one applies
$0\to\mathcal O\to\mathcal O(E_\beta)
\to\mathcal O_{E_\beta}(-1)\to0$ to the usual
blowup structure-sheaf calculation. Projection formula,
the divisor ideal sequence and the absence of second
direct images then give
\[
 \beta_*\mathcal O_{D'}=\mathcal O_{D''},
 \qquad R^1\beta_*\mathcal O_{D'}=0.                         \tag{P4R.13.2}
\]
These particular maps identify $H^0(D,\mathcal O_D)$
with $H^0(D_0,\mathcal O_{D_0})$, where
\[
 D_0=2H_0+V_0,
 \qquad V_0=\sum_j a_jF_j
\]
on the minimal ruled surface. The kernel of restriction
from $D_0$ to $2H_0$ is $\mathcal O_{V_0}(-2H_0)$.
Its successive layers on each thick fibre $a_jF_j$ are
$\mathcal O_{\mathbb P^1}(-2)$, since $F_j^2=0$ and
$H_0F_j=1$. They have no global sections. Restriction
to $2H_0$ is therefore injective on global functions.
This double section is Cohen--Macaulay and has no
zero-dimensional associated point, so a section which
vanishes generically on it vanishes everywhere.

\item Let $C_{E,\rm sch}$ be the primary component of
$C_{\rm sch}$ supported on $C_E$: locally its ideal
is obtained by localizing the conductor ideal at the
generic point of $C_E$ and contracting. Primary
decomposition of the pure one-dimensional conductor
quotient makes these ideals coherent and compatible
on overlaps. Restriction is a surjection
\[
 \mathcal O_{C_{\rm sch}}\longrightarrow
 \mathcal O_{C_{E,\rm sch}}.
\]
Its map on $H^0$ is injective. Indeed, a global element
of its kernel maps by (P4R.13.1) to a function on $D$
which vanishes generically on $2H$; (P4R.13.2) and the
preceding paragraph make it zero. Its map on $H^1$
is surjective, because the kernel is supported in
dimension at most one and has no second cohomology.
We have proved
\[
 \chi(\mathcal O_{C_{\rm sch}})
 \le\chi(\mathcal O_{C_{E,\rm sch}}).                       \tag{P4R.13.3}
\]

\item We evaluate the right side using the actual primary
ideal, without removing the rational components from
the left side. In type $A_2$ the primary conductor
is reduced: a pure one-dimensional quotient which
is generically reduced has no nilpotent submodule
of finite support. Thus its Euler characteristic
is $-\sum_p\delta(C_E,p)\le0$. In type $D_4$ the
primary ideal is the symbolic square of the ideal
of $C_E$, by exactly the generic-conductor and
depth-one argument of P4.31. P4.30 gives the same
nonpositive bound, with equality only for smooth
$C_E$.

\item In type $A_3$, the primary ideal contains that
symbolic square and its nilpotent module is
torsionfree of rank one on $C_E$. The proof of
P4R.10 applies to this primary quotient: its
differentiation map, saturated Hessian line and
projection to the quotient conormal line depend
only on this ideal and on the actual equation of
$S$. The presence of other conductor components
changes none of those maps. More explicitly its
Euler characteristic is
\[
 -\sum_p\delta(C_E,p)+\deg R
 -\frac12\sum_p d_p-\sum_p b_p-\sum_p q_p\le0,
\]
with the same ramification divisor $R$, Hessian
zero orders $d_p$, projected lattice colengths
$b_p$ and descent defects $q_p$ as in P4R.10.
The local inequalities there prove this bound
at every singular point of $C_E$. Equality
makes $C_E$ smooth and its Hessian nowhere zero.

\item On the other hand P4.33 gives
$\chi(\mathcal O_{C_{\rm sch}})=g-1$. Thus
\[
 0\le g-1
 =\chi(\mathcal O_{C_{\rm sch}})
 \le\chi(\mathcal O_{C_{E,\rm sch}})\le0.                    \tag{P4R.13.4}
\]
All these quantities vanish. We obtain $g=1$,
$C_E=E$ smooth, and an unramified map $H\to E$.
In type $A_3$ its transverse Hessian has rank
one everywhere on $E$.

\item It remains to exclude attachments to this smooth
elliptic component. Put
\[
 K=\mathcal I_E/\mathcal I_E^2,
 \qquad N_E=\mathcal O_Y(S)|_E.
\]
Both $\deg K$ and $\deg N_E$ are zero. The first
follows from the conormal exact sequence,
$c_1(TY)E=0$, and $\deg\Omega_E=0$; the second
follows because a positive multiple of $S$ is
principal in a tube of this fibre.

\item In type $A_2$ let $L\subset K$ be the saturated
quadratic-image line and put $Q=K/L$. The nonzero
quadratic tensor is a section of
$L^2\otimes N_E$, so $2\deg L\ge0$.
The cubic jet on the null direction is an
intrinsic section of $Q^3\otimes N_E$.
Here is the coordinate check, including points
where the quadratic tensor vanishes. A change
of a second-order lift changes the cubic
coefficient by the polarization of the
quadratic tensor. That polarization vanishes
on the null direction. A change of tangential
lift contributes a tangential derivative of
the rank-one quadratic tensor; its evaluation
twice on the null direction is again zero.
Changing a reduced defining equation multiplies
the tensor by its transition unit. Thus the
stated line bundle and its holomorphic section
are defined on all of $E$. The section is
nonzero generically by type $A_2$, so
$-3\deg L\ge0$. Both degrees are zero,
and both sections are nowhere zero. The
parametrized Morse lemma and a cube root
of the residual unit now give the product
equation $X^2+Y^3=0$ at every point of $E$.

\item In type $A_3$ (P4R.13.4) already makes the
quadratic tensor nowhere zero, whence
$\deg L=\deg Q=0$. Its cubic null jet is
identically zero. Removing the quadratic
direction gives an intrinsic residual fourth
jet in $Q^4\otimes N_E$. In local coordinates
its relevant terms and coefficient are
\[
 a x^2+bxy^2+cy^4,
 \qquad q_4=c-\frac{b^2}{4a},\qquad a\ne0.                 \tag{P4R.13.5}
\]
Substitution of the critical coordinate
$x=-by^2/(2a)+O(y^3)$ proves the formula.
The change $x\mapsto x+vy^2$ sends $b$ to
$b+2av$ and $c$ to $c+bv+av^2$, leaving
$q_4$ unchanged. Changes of normal frames
and of equation give the claimed
$Q^4\otimes N_E$ transition. Tangential
lift changes do not affect a function
whose residual order is four. This section
is nonzero generically and has degree zero;
it never vanishes. The residual function
is $y^4$ times a unit, giving a product
tacnode at every point of $E$.

\item In type $D_4$, the equation belongs to
$\mathcal I_E^3$. This ideal is primary
because $E$ is smooth, so generic order
three proves membership everywhere.
The cubic normal jet belongs to
$\operatorname{Sym}^3K\otimes N_E$.
Its discriminant is a section of
$(\det K)^6\otimes N_E^4$, of degree
zero. It is nonzero generically, hence
nowhere zero. The three distinct normal
tangent directions give three smooth
holomorphic branches locally along $E$.
After straightening the first two to
$x=0,y=0$, the third is $ax+by=0$
with both coefficients units. Setting
$X=ax,Y=-by$ gives the product equation
$XY(X-Y)=0$.

\item In all three cases the entire singular
locus near every point of $E$ is $E$
itself. A rational conductor tree could
not meet it. Connectedness of $C$ then
forces $C=E$. The displayed local
product forms make $N$ smooth along
its conductor, and P4R.6 already makes
it smooth elsewhere. Hence $T=N$.
A finite normalization cannot contract
a curve, and a rational curve cannot
map nonconstantly to $E$. Thus $D=2H$
has no vertical component. The first
blowup of $2H_0$ in a nonempty actual
blowup sequence would create a vertical
pole component of coefficient one;
later blowups never remove its strict
transform from the pole support.
Consequently no blowup occurs and $N$
is minimal elliptic ruled. Finally
\[
 e(S)=e(N)-e(H)+e(E)=0.
\]
The exact coefficient changes, discriminant
weights and degree identities are executed
in \texttt{code/part4/elliptic\_primary\_restriction.py};
its JSON and text outputs are retained.

\end{enumerate}
\end{proof}
